\documentclass[11pt]{article}
\usepackage{amsmath,amssymb}
\usepackage[margin=1in]{geometry}

\title{Step 275: QUE Classification Within the Framework}
\author{Six Birds Foundations III}
\date{}

\begin{document}
\maketitle

\section*{Carrier}
Let \(\Gamma=\mathrm{SL}_2(\mathbb Z)\) or an arithmetic congruence subgroup, and let
\(\{\phi_j\}\) be an orthonormal Hecke--Maass cusp-form basis on
\(\Gamma\backslash\mathbb H\), with Laplace eigenvalue \(1/4+t_j^2\).  For
\(f\in C(\Gamma\backslash\mathbb H)\), define
\[
\Xi_{\mathrm{QUE}}(j;f)=
\left|\int_{\Gamma\backslash\mathbb H} f(z)|\phi_j(z)|^2\,d\mathrm{vol}(z)
-\frac{1}{\mathrm{vol}(\Gamma\backslash\mathbb H)}
\int_{\Gamma\backslash\mathbb H}f(z)\,d\mathrm{vol}(z)\right|.
\]
QUE closure is \(\Xi_{\mathrm{QUE}}(j;f)\to 0\) for every continuous \(f\) as
\(t_j\to\infty\).

\section*{Framework Classification}
The primary classification is the Selberg-Class Subconvexity Extension, Type
\(\alpha\), via Watson's triple-product bridge.  Quantitative QUE periods such as
\(\int \phi_j^2\psi\) are related to central triple-product \(L\)-values, in
particular central values of the shape
\[
L\!\left(\frac12,\operatorname{sym}^2\phi_j\times\psi\right),
\]
up to the usual local factors and normalization.  Thus the arithmetic effective
QUE route is a central-\(L\)-value size route.

The secondary classification is Cross-Correlation Extension Type Ia, since the
QUE residual is also a mass/eigenfunction correlation against a test observable.
This secondary shadow does not displace the primary subconvexity classification.

\section*{Exclusions}
QUE is not an SCDG instance: it is not a per-\(L\)-function RH-analog or
zero-location defect.  No new 11th framework finding is required.

\section*{Verdict}
\[
\boxed{\texttt{V\_que\_in\_subconvexity}}
\]
Arithmetic QUE in the cited Hecke settings is external known content; general
non-arithmetic QUE remains open.  No RH or GRH claim is made.

\end{document}
